\chapter*{Tercera nota de actualización, 12 de septiembre de 2026}
\addcontentsline{toc}{chapter}{Tercera nota de actualización, 12 de septiembre}

\titular{\textbf{La Ruta A existía}: los analíticos del proyecto estaban publicados bajo un
prefijo de dirección distinto del que este documento sondeó tres días. Con ellos, \textbf{el
perímetro pensionario se reparte por institución} ---el IMSS aporta 125.702,8 de los
136.586,9 mdp del aumento---, la salud del IMSS se separa de sus pensiones, y la auditoría de
etiquetado se hace. Y una corrección que no dependía de la ruta: \textbf{la caída del gasto
corriente estructural no se explica sólo por la exclusión de las empresas públicas}.}

\textbf{Qué es esta nota y qué no es.} Es la tercera adenda y obedece a la misma regla:
\textbf{el cuerpo de los dieciséis capítulos no se reescribe}. A diferencia de las dos
anteriores, \textbf{esta nota corrige afirmaciones del cuerpo}, y dice cuáles. Ninguna cifra
del cuerpo estaba mal calculada; lo que cambia es lo que el cuerpo dijo que \emph{no se podía
saber} y dos lecturas que excedían su cálculo. Lo que se añade sale de fuentes que se
leyeron el 12 de septiembre; el documento de CIEP sobre el mismo paquete, publicado el 11, se
leyó ese mismo día, \textbf{después} de cerrar las dos notas anteriores.

\section*{Uno. La ruta, y lo que el registro dijo mal}

El registro y las dos notas afirman que los cuatro analíticos del proyecto respondían 404
porque \textbf{«el árbol de 2027 no existe»}, y que la exposición de motivos del PPEF no se
servía \textbf{por sexto ejercicio consecutivo}. \textbf{Las dos afirmaciones son falsas
para 2027.} Los 404 eran ciertos; el diagnóstico, no: se sondeó la ruta de los ejercicios
anteriores. La Secretaría publicó el paquete bajo el prefijo
\texttt{ppef.hacienda.gob.mx/work/models/PP3F2709/PPEF2027/yik327fP/}, que este proyecto no
conocía y que localizó en las referencias del documento de CIEP.

\begin{itemize}
\item \textbf{La Carta del proyecto} está ahí con fecha de servidor del \textbf{8 de
septiembre}, 08:58 GMT: existía el día de la entrega. La exposición de motivos también se
sirve: se verificó el documento completo y su primer capítulo.
\item \textbf{Los cuatro analíticos} se sirven en la carpeta de analíticos presupuestarios de
ese árbol, \textbf{sin} el segmento «Proyecto» ---con él siguen dando 404---, con fecha de
servidor del 11 de septiembre, 20:55 GMT. Esa fecha es posterior a la tercera sonda de este
documento; \textbf{no se sabe si es la de primera publicación o la de una resubida}, y no se
afirma ninguna de las dos.
\end{itemize}

\textbf{Cambio de ruta: 12 de septiembre, 23:08 --- Ruta A.} Los analíticos se descargaron,
y antes de usarlos se probaron: los cuarenta ramos del Anexo 1 salen del analítico en los
dos ejercicios con una diferencia máxima de 0,05 mdp, y el IMSS, el ISSSTE, Pemex y CFE de
2027 salen al décimo.

\porqueimporta{\textbf{Un 404 sobre una ruta conocida no prueba que el recurso no exista.}
Este documento distinguió con cuidado entre «árbol no creado», para los analíticos, y
«servicio caído», para el portal de datos abiertos; \textbf{el primero de los dos diagnósticos
era de ruta}. El segundo sigue en pie: esas ranuras responden 403 el 12 de septiembre. La
regla que queda: antes de declarar
que una fuente no se publicó, se busca su dirección vigente en lo que sí se publicó.}

\textbf{El diff institucional}, que la Ruta B no permitió, se corrió: \textbf{ningún ramo
nace ni se extingue} entre 2026 y 2027, y casi todo lo que aparece como programa nuevo es
recodificación ---el Ramo 28 abre en fondos con nombre propio lo que era un solo concepto;
CFE renumera tres programas; el programa de atención gratuita de IMSS-Bienestar se divide en
tres---. \textbf{La Guardia Nacional sigue en el Ramo 07 y crece 28,9~\% real}, sin cambio de
ramo; el Cuerpo de Policía Militar desaparece como unidad responsable, y el analítico no
permite atribuir a dónde fue su presupuesto.

\section*{Dos. La regla fiscal: por qué cae el gasto corriente estructural}

El capítulo de la regla fiscal y la implicación 2 atribuyen la caída del agregado ---de
3.868.319,9 a 1.884.958,2 mdp--- a que la reforma de abril a la LFPRH dispone que el gasto
de las empresas públicas «no se contabilizará». \textbf{Es incompleto, y la fuente que lo
completa estaba en carpeta desde el 8 de septiembre}: el propio CGPE publica la metodología
nueva en su página 24.

\begin{table}[htbp]\centering\footnotesize
\caption{El límite máximo del gasto corriente estructural con la metodología nueva (mdp)}\label{cua:gcemetodo}
\begin{tabular}{>{\raggedright\arraybackslash}p{9.6cm}r}
\toprule
Concepto & Monto \\
\midrule
\textbf{A) Gasto neto pagado CP 2025} & \textbf{8.374.179,5} \\
\quad (1) Costo financiero & 1.131.312,2 \\
\quad (2) Participaciones & 1.350.864,7 \\
\quad (3) Adefas & 28.863,6 \\
\quad (4) Pensiones y jubilaciones & 1.458.485,7 \\
\quad (5) Inversión física y financiera directa GF y ECPD & 643.527,6 \\
\quad (6) Inversión física y financiera directa ECPI & 35.936,0 \\
\quad \textbf{(7) Programas sociales universales establecidos en la CPEUM} & \textbf{766.412,4} \\
\quad \textbf{(8) Servicios personales función educación} & \textbf{648.331,6} \\
\quad \textbf{(9) Servicios personales función salud} & \textbf{522.120,5} \\
\quad \textbf{(10) Servicios personales función asuntos de orden público y de seguridad interior} & \textbf{27.635,3} \\
\textbf{B) Gasto corriente estructural pagado CP 2025} & \textbf{1.760.689,7} \\
\quad Diferimiento de pagos & 75.720,6 \\
\textbf{C) Gasto corriente estructural devengado CP 2025} & \textbf{1.836.410,3} \\
\textbf{D) Límite máximo para 2027} & \textbf{2.058.457,0} \\
\midrule \textbf{Suma de los rubros (7) a (10)} & \textbf{1.964.499,8} \\
\bottomrule
\end{tabular}
\\[2pt]\begin{minipage}{0.92\textwidth}\scriptsize Fuente: CGPE 2027, edición SHCP, p. 24, extraído del PDF a \texttt{datos/gce\_limite\_2027.csv}. El renglón A excluye, además, el gasto de las empresas públicas del Estado (nota 2/ del propio cuadro, artículo 2, fracción XXIV Bis, de la LFPRH). \textbf{Los rubros en negritas no se restaban antes}; son de la Cuenta Pública 2025, no del proyecto 2027. La D resulta de actualizar C por precios (3,8 y 4,0~\%) y crecimiento real (1,9~\% cada año).\end{minipage}
\end{table}


El cómputo nuevo resta, además de las empresas públicas, \textbf{cuatro rubros que antes no
restaba}: los programas sociales universales establecidos en la Constitución y los servicios
personales de las funciones de educación, salud y orden público. Sobre la Cuenta Pública
2025, esos cuatro rubros suman \textbf{1.964.499,8 mdp}.

\porqueimporta{\textbf{La frase correcta es otra}: el agregado que la regla mide ya no excluye
sólo a las empresas públicas, sino también a los programas sociales constitucionales y a la
nómina de tres funciones. La cifra de los rubros nuevos es de la Cuenta Pública 2025 y la
caída del agregado es entre dos proyectos, de modo que \textbf{no se restan una de otra}; lo
que se puede decir es que son del mismo orden de magnitud, y que el gasto corriente de las
empresas públicas sin costo financiero ---573.721,3 mdp en el proyecto 2027, según los
analíticos--- es bastante menor.}

\textbf{El punto que el capítulo dejó «ni se afirma ni se niega» queda resuelto.} Sí hay
aportación de capital a una empresa pública clasificada en el capítulo 7000: la partida
\textbf{73903} del Ramo 18 lleva \textbf{81.103,0 mdp} en 2027, contra 263.476,3 en 2026. Con
la partida exacta, la aritmética del capítulo de energía se afina: el resto del Ramo 18 sube
\textbf{568,4 mdp}, no 571,1, que salía de las cifras redondeadas de la prosa del CGPE.

\section*{Tres. Pensiones por institución, y dos advertencias al cuerpo}

\begin{table}[htbp]\centering\footnotesize
\caption{El perímetro pensionario por institución, y el Ramo 19 bruto y neto, línea P (mdp)}\label{cua:pensinst}
\begin{tabular}{>{\raggedright\arraybackslash}p{4.6cm}rrrr}
\toprule
Institución & 2026 & 2027 & Dif. nominal & Var. real (\%) \\
\midrule
IMSS & 1.010.525,2 & 1.136.228,0 & 125.702,8 & +8,1 \\
ISSSTE & 395.156,0 & 409.867,0 & 14.711,0 & $-$0,3 \\
Pemex & 92.333,4 & 94.472,9 & 2.139,5 & $-$1,6 \\
CFE & 67.841,7 & 66.154,8 & $-$1.686,9 & $-$6,2 \\
Resto del Gobierno Federal, por diferencia & 138.302,9 & 134.023,4 & $-$4.279,5 & $-$6,8 \\
\midrule \textbf{Perímetro del Anexo 3} & \textbf{1.704.159,2} & \textbf{1.840.746,1} & \textbf{136.586,9} & \textbf{+3,9} \\
\midrule \multicolumn{5}{l}{\emph{Ramo 19: no se suma a lo anterior}} \\
Ramo 19 bruto & 1.541.518,7 & 1.661.007,9 & 119.489,2 & +3,6 \\
Transferido a IMSS e ISSSTE (UR GYR y GYN) & 1.401.952,8 & 1.526.114,6 & 124.161,8 & +4,7 \\
Ramo 19 neto de esas transferencias & 139.565,9 & 134.893,2 & $-$4.672,7 & $-$7,1 \\
\bottomrule
\end{tabular}
\\[2pt]\begin{minipage}{0.92\textwidth}\scriptsize Fuente: analíticos del proyecto 2026 y 2027, tipo de gasto 4 (pensiones y jubilaciones), \texttt{datos/ruta\_a/pensiones\_tg4\_institucion.csv}; perímetro por diferencia de los dos renglones del Anexo 3 de los decretos. Deflactor del PIB 1,040. \textbf{El Ramo 19 bruto cuenta dos veces lo que transfiere al IMSS y al ISSSTE}, que reaparece en el presupuesto de esas entidades; el neto es la cifra que publica CIEP.\end{minipage}
\end{table}


\textbf{El perímetro que este documento adjudicó por diferencia del Anexo 3 se descompone
ahora sin residuo inexplicado.} El IMSS pasa de 1.010.525,2 a \textbf{1.136.228,0 mdp},
$+$8,1~\% real, y aporta el 92~\% del aumento del perímetro; el ISSSTE, Pemex, CFE y el resto
del Gobierno Federal caen en términos reales. El «qué no sabemos» del capítulo de pensiones
queda respondido.

\textbf{Primera advertencia: el Ramo 19 del cuadro de gasto total.} El cuadro de los ramos
que más se mueven pone al Ramo 19 en segundo lugar, con $+$119.489,2 mdp. La cifra es
correcta y es bruta: \textbf{lo que el Ramo 19 transfiere al IMSS y al ISSSTE crece
124.161,8 mdp}, más que todo el aumento del ramo, y esos mismos pesos reaparecen en el
presupuesto del IMSS. Neto de esas transferencias, el Ramo 19 \textbf{cae}. Leer el renglón
bruto como una prioridad de gasto nueva sería contar dos veces.

\textbf{Segunda advertencia: el titular de pensiones por persona.} El resumen y la
implicación 4 dicen que, por persona de 65 años y más, el gasto pensionario \textbf{cae
0,4~\% real}. Con el deflactor del PIB (1,040) es así. Con el índice que usa el propio CGPE
en sus cuadros (1,0325), el mismo cociente \textbf{sube 0,3~\%}. \textbf{Una cifra que cambia
de signo con la convención de deflactación no sostiene un verbo}: lo que el documento puede
afirmar es que el perímetro pensionario \textbf{crece prácticamente al mismo ritmo} que la
población de 65 años y más, no que caiga por persona. La nota de método promete una línea de
sensibilidad para estos casos, y este titular no la tuvo.

\section*{Cuatro. Salud del IMSS y del ISSSTE, separada de sus pensiones}

\begin{table}[htbp]\centering\footnotesize
\caption{IMSS e ISSSTE por función, y la función salud completa, línea P (mdp)}\label{cua:imssfun}
\begin{tabular}{ll>{\raggedleft\arraybackslash}p{2.0cm}>{\raggedleft\arraybackslash}p{2.0cm}r}
\toprule
Entidad & Función & 2026 & 2027 & Var. real (\%) \\
\midrule
IMSS & 2.3 salud & 556.342,6 & 640.797,1 & +10,8 \\
IMSS & 2.6 protección social & 1.032.733,4 & 1.160.838,9 & +8,1 \\
ISSSTE & 2.3 salud & 81.525,8 & 99.963,9 & +17,9 \\
ISSSTE & 2.6 protección social & 457.005,5 & 471.135,6 & $-$0,9 \\
\midrule \textbf{Función salud, GF y entidades} & 2.3 & \textbf{974.302,2} & \textbf{1.114.506,5} & \textbf{+10,0} \\
\bottomrule
\end{tabular}
\\[2pt]\begin{minipage}{0.92\textwidth}\scriptsize Fuente: analíticos del proyecto por ramo, función, UR y objeto del gasto, de Gobierno Federal y de entidades; \texttt{datos/ruta\_a/entidades\_funcion.csv} y \texttt{salud\_subfuncion.csv}. La función 2.6 del IMSS y del ISSSTE contiene sus pensiones. La función salud de 2026 reproduce la del documento 2026, 974.302,2.\end{minipage}
\end{table}


El capítulo de salud tituló que el presupuesto del IMSS crece 9,0~\% y advirtió, con razón,
que esa cifra mezcla atención médica y pensiones. \textbf{Separadas, la función salud del
IMSS crece 10,8~\% real y la del ISSSTE 17,9~\%}; la protección social del ISSSTE, que
contiene sus pensiones, cae. La función salud completa, del Gobierno Federal y las
entidades, pasa de 974.302,2 a \textbf{1.114.506,5 mdp}, $+$10,0~\% real.

Una caída que \textbf{no debe leerse}: la subfunción 2.3.4 cae 74,3~\%, y casi toda es el
Ramo 33 ---el FASSA--- que baja en esa subfunción lo que sube en la 2.3.2. Es el mismo fondo
reclasificado entre subfunciones, no un recorte.

\section*{Cinco. Anexos transversales: la auditoría que faltaba, y el anexo que faltaba}

\textbf{El capítulo de anexos dice «aparece un anexo nuevo, el 33». Son dos.} El proyecto de
decreto trae también el \textbf{Anexo 32, Acciones para la Movilidad y Seguridad Vial},
con \textbf{175.045,5 mdp}, y el artículo 3 del propio decreto los nombra a los dos. El 32 no
se leyó porque en el decreto \textbf{no tiene encabezado ni renglón de total}: su tabla
empieza a media página, y el procedimiento de lectura reconocía cada anexo por esos dos
renglones. Con él, la suma aritmética de los anexos del cuadro \ref{cua:trans} sería
de 4.600.786,6 mdp; sigue sin medir ningún gasto.

Las ranuras de datos abiertos siguen sin servir el etiquetado por programa, así que la
auditoría se hizo \textbf{leyendo las tablas de los anexos del decreto}, que desde 2027 traen
objetivo, ramo, programa y acción, y validando cada extracción contra el total de su anexo.

\begin{table}[htbp]\centering\footnotesize
\caption{La auditoría de etiquetado de los anexos transversales}\label{cua:auditoria}
\begin{tabular}{>{\raggedright\arraybackslash}p{5.0cm}r>{\raggedleft\arraybackslash}p{2.2cm}>{\raggedleft\arraybackslash}p{2.8cm}}
\toprule
Indicador & 2026 & 2027, mismos anexos & 2027, con el 32 y el 33 \\
\midrule
Programas etiquetados & 226 & 228 & 243 \\
Programas en más de un anexo & 123 & 123 & 136 \\
\textbf{\% del etiquetado en programas con más de una etiqueta} & \textbf{95,4} & \textbf{95,8} & \textbf{95,2} \\
Suma aritmética de etiquetas (mdp) & 4.284.445,1 & 4.381.269,1 & 4.599.683,9 \\
\textbf{Pensiones no contributivas en el Anexo 13 (\%)} & \textbf{48,1} & \textbf{57,2} & \textbf{57,2} \\
\midrule Anexo 32: programas ya etiquetados en otro anexo & --- & --- & 16 de 23 \\
Anexo 32: mdp ya etiquetados en otro anexo & --- & --- & 130.912,2 de 175.045,5 \\
\midrule Anexo 33: programas ya etiquetados en otro anexo & --- & --- & 54 de 62 \\
Anexo 33: mdp ya etiquetados en otro anexo & --- & --- & 43.322,2 de 43.369,4 \\
\bottomrule
\end{tabular}
\\[2pt]\begin{minipage}{0.92\textwidth}\scriptsize Fuente: tablas de los anexos 13 a 19, 31, 32 y 33 de los proyectos de decreto, extraídas programa por programa y validadas contra el «Total general» de cada anexo (cobertura del 100~\% salvo el Anexo 13 de 2027, 99,99~\%, y el 31, 99,78~\%; el 32 no imprime total). \texttt{datos/ruta\_a/anexos\_auditoria.csv}. Pensiones no contributivas: claves S176, S286 y S316. \textbf{La suma de etiquetas no mide gasto.}\end{minipage}
\end{table}


\textbf{Las tres preguntas que el capítulo dejó sin responder, respondidas:}
\begin{itemize}
\item \textbf{El traslape no cambia:} el 95,8~\% de lo etiquetado en los anexos que existían
en 2026 está en programas que llevan más de una etiqueta; en 2026 era el 95,4~\% con el mismo
procedimiento.
\item \textbf{El anexo de igualdad se vuelve más pensión:} el 57,2~\% son tres pensiones no
contributivas, contra el 48,1~\% de 2026.
\item \textbf{Los dos anexos nuevos son, sobre todo, etiqueta sobre programas ya
etiquetados}: el 99,9~\% de los pesos del Anexo 33 y el 74,8~\% de los del Anexo 32 están
en programas que ya llevaban otra etiqueta en 2027. Responde el «qué no sabemos» del
capítulo sobre el Anexo 33.
\end{itemize}

\begin{table}[htbp]\centering\footnotesize
\caption{Los cinco programas que más etiqueta pierden en el Anexo 16, cambio climático (mdp)}\label{cua:anexo16}
\begin{tabular}{l>{\raggedright\arraybackslash}p{4.6cm}rrr}
\toprule
Ramo-programa & Nombre & Etiqueta 2026 & Etiqueta 2027 & Presupuesto 2026 \\
\midrule
07-A001 & Defensa de la Integridad, la Independencia, la Soberanía del Territorio Nacional & 44.673,4 & 0,0 & 44.673,4 \\
07-K003 & Infraestructura en materia de seguridad nacional & 40.000,0 & 0,0 & 1.510,9 \\
18-P038 & Articulación de la política de hidrocarburos & 26.370,0 & 0,0 & 263.700,0 \\
07-E017 & Servicios de dirección la infraestructura aeroportuaria, ferroviaria y de servicios auxiliares & 3.422,2 & 21,0 & 3.422,2 \\
07-E015 & Servicios Públicos de Transporte Masivo de Personas y Carga Tren Maya & 744,1 & 139,4 & 744,1 \\
\midrule \textbf{Anexo 16 completo} & & \textbf{212.569,7} & \textbf{160.120,9} &  \\
\bottomrule
\end{tabular}
\\[2pt]\begin{minipage}{0.92\textwidth}\scriptsize Fuente: Anexo 16 de los proyectos de decreto 2026 y 2027; presupuesto del programa, del analítico del proyecto 2026. \texttt{datos/ruta\_a/anexo16\_programas.csv}. \textbf{La etiqueta de 2026 del K003 es más de 26 veces el presupuesto del programa}, y el programa no existe en 2027.\end{minipage}
\end{table}


\textbf{La caída del Anexo 16 es de etiquetado.} Tres programas dejan de etiquetarse y suman
111.043,4 mdp: dos de Defensa y el de articulación de la política de hidrocarburos, cuyo
presupuesto se redujo por la reformulación del rescate de Pemex. \textbf{No es la conclusión
de la inversión del Tren Maya}: su programa de servicios de transporte pierde 604,8 mdp de
etiqueta. Y una anomalía de la fuente, que se registra sin interpretarla: \textbf{en 2026 el
anexo etiquetó 40.000,0 mdp a un programa cuyo presupuesto en el analítico del proyecto era
de 1.510,9}.

\section*{Seis. Lo que estaba en el CGPE y en los Pre-Criterios, y no se usó}

Tres omisiones propias que \textbf{no dependían de la ruta}, y que el documento declaró como
límites de la fuente:
\begin{itemize}
\item \textbf{La inversión por proyecto.} El capítulo de inversión escribió que la
composición por proyecto «exige el analítico del proyecto y la cartera». La página 32 del
CGPE publica las prioridades de inversión por proyecto, con un total de \textbf{560.172,6
mdp}.
\item \textbf{La pensión no contributiva por programa.} El capítulo de pensiones escribió que
«se puede acotar con el Ramo 20, pero no separar». La misma página publica las tres pensiones
no contributivas, que suman \textbf{640.291,0 mdp}.
\item \textbf{La meta fiscal de abril.} El capítulo macro comparó los Pre-Criterios con el
CGPE sólo en variables macroeconómicas. \textbf{Los Pre-Criterios fijaban unos RFSP de
3,5~\% del PIB para 2027}; el CGPE los deja en 3,9. El saldo proyectado, 55,0~\% del PIB, es el
mismo porcentaje en los dos documentos sobre dos PIB distintos.
\end{itemize}

\section*{Siete. Qué se volvió a correr}

\begin{itemize}
\item \textbf{Identidades del paquete: de 49 a 52 pruebas, ninguna falla.} Las tres nuevas
corren solas al estar los analíticos de 2027 en carpeta: el bruto del Anexo 1 contra los
analíticos del Gobierno Federal y de las entidades, y el corte por función contra el corte
por programa en cada uno de los dos, las tres al peso.
\item \textbf{Guion de la Ruta A}, nuevo: los cuarenta ramos y las cuatro entidades contra el
Anexo 1, la función salud y la partida 73903 de 2026 contra el documento de ese año, el Ramo
19 neto, el Anexo 32, el cuadro del gasto corriente estructural contra su propia aritmética
y las tablas de la página 32 del CGPE: \textbf{catorce pruebas, ninguna falla}.
\item \textbf{Compuerta semántica}, aplicada al titular, a los cinco pies de cuadro y a las
frases de «por qué importa» de esta nota; el detalle está en la bitácora.
\item \textbf{Checklist de omisiones}: B2 y B3 quedan resueltas a medias ---la meta de abril
sí, el pronóstico externo no---; B5 y B10 resueltas; C5 encuentra su fuente en la página 32 del
CGPE, de la que esta nota sólo cita el total; \textbf{B1, la vara de suficiencia en salud, sigue fuera, ahora por cuarto ejercicio}
contando el de 2025.
\end{itemize}

\textbf{Lo que sigue sin hacerse}, aunque la Ruta A ya lo permitiría: el agua dentro del Ramo
16; la lectura de la exposición de motivos del PPEF; la línea G por programa; y el cruce de
la ZAP con el gasto, que el analítico no permite porque llega a entidad federativa y no a
municipio.

\begin{ficha}{tercera nota de actualización}
\fila{Perímetro}{Los del cuerpo, más: tipo de gasto 4 de los analíticos por institución;
funciones 2.3 y 2.6 de las entidades; partida 73903; tablas de los anexos transversales
programa por programa; y el cuadro de la página 24 del CGPE, que es de Cuenta Pública 2025.}
\fila{Línea de comparación}{Línea P en todos los cuadros de los analíticos. El cuadro del gasto
corriente estructural no tiene línea: es un cómputo de la Cuenta Pública 2025 proyectado.}
\fila{Deflactor}{Del PIB, 1,040. La sensibilidad al índice de 1,0325 se reporta sólo donde
cambia el signo de una afirmación del cuerpo.}
\fila{Año de los pesos}{Pesos corrientes en los niveles; pesos de 2027 en las variaciones.
Pesos de 2025 en el cuadro \ref{cua:gcemetodo}.}
\fila{Denominadores}{Población de 65 años y más, de la capa demográfica, sólo en la
sensibilidad del titular de pensiones.}
\fila{Fuentes}{Los cuatro analíticos del proyecto 2027 y los de 2026, en la carpeta de cada
ejercicio con su sha256 en el manifiesto; los proyectos de decreto 2026 y 2027; CGPE 2027,
pp. 24 y 32; Pre-Criterios 2027, p. 32. Los CSV de esta nota están en
\texttt{datos/ruta\_a/} y los genera \texttt{ruta\_a.py}, que viaja en la carpeta del
proyecto y no en este paquete porque exige los analíticos.}
\fila{Qué no puede afirmarse con ellas}{Que la caída del gasto corriente estructural se
descomponga exactamente en los cuatro rubros nuevos: son de otro año y de otra base. Que el
Anexo 33 no tenga pesos nuevos: el 0,1~\% restante son programas que no llevaban otra
etiqueta, no necesariamente gasto nuevo. Ni que la fecha del servidor de los analíticos sea
la de su primera publicación.}
\end{ficha}
